\documentclass[11pt]{article}

\usepackage[margin=1in]{geometry}
\usepackage{amsmath,amssymb}
\usepackage{xcolor}
\usepackage[most]{tcolorbox}
\usepackage{fancyvrb}
\usepackage{enumitem}
\usepackage{parskip}
\usepackage{hyperref}
\hypersetup{colorlinks=true,linkcolor=blue!50!black,urlcolor=blue!50!black}

\definecolor{ideabg}{RGB}{255,249,230}
\definecolor{ideaframe}{RGB}{200,120,0}
\definecolor{codebg}{RGB}{245,245,245}

\newtcolorbox{keyidea}{
  colback=ideabg, colframe=ideaframe, boxrule=0.6pt, arc=2pt,
  left=8pt, right=8pt, top=7pt, bottom=7pt, breakable
}

\newcommand{\code}[1]{\texttt{#1}}

\title{\vspace{-1.5em}How the Graded DSL Works\\[0.3em]
  \large An embedded, type-level-graded language for safe floating point}
\author{}
\date{\today}

\begin{document}
\maketitle
\vspace{-2em}

\noindent
This note explains, from the ground up, how the little language in
\code{src/graded.rs} works. It assumes no background in programming-language
theory. One running example is used throughout: the formula

\begin{center}
\code{(var(0) * var(0) + var(1) * var(1)).sqrt()}
\end{center}

\noindent
which is just the length of a vector, $\sqrt{x^{2}+y^{2}}$. Four questions are
answered in turn: what an \emph{embedded DSL} is, what the \emph{type checker}
did, in what sense the language is \emph{interpreted}, and the role of
\emph{operator overloading}.

\section{What an embedded DSL is}

A \textbf{DSL} (domain-specific language) is a small language built for one
narrow job. SQL is one: its job is database queries. Regular expressions are
one: their job is text patterns.

There are two ways to make a DSL:

\begin{itemize}[leftmargin=1.4em]
  \item \textbf{External DSL.} You invent brand-new syntax and write a program
  (a \emph{parser}) that reads that text and makes sense of it. \code{SELECT *
  FROM users} is not valid in any ordinary programming language; somebody wrote
  a parser that understands those words.
  \item \textbf{Embedded DSL.} You do \emph{not} invent new syntax. You write
  your little language using the ordinary grammar of a \emph{host} language ---
  here, Rust --- and simply arrange for the pieces to \emph{mean} something
  special. It is ``embedded'' because it lives inside Rust; the Rust compiler
  reads it as ordinary Rust.
\end{itemize}

Our language is \textbf{embedded}. When you write \code{var(0) * var(0)}, you
are writing one hundred percent legal Rust. There is no parser and no new
grammar. The ordinary Rust pieces have just been given a special meaning.

Here is the one idea that unlocks everything else.

\begin{keyidea}
\textbf{When you write \code{var(0) * var(0)}, no multiplication happens.} It
\emph{looks} like arithmetic, but nothing is computed. Instead, Rust builds a
little model --- think of it as writing the recipe step \emph{``multiply these
two things''} on an index card, rather than actually multiplying. The whole
formula becomes a stack of index cards describing what to do. The actual numbers
do not exist yet.
\end{keyidea}

So \code{(var(0)*var(0) + var(1)*var(1)).sqrt()} does not produce a number. It
produces a little \textbf{tree of index cards}:

\begin{center}
\begin{BVerbatim}
            sqrt
             |
            add
          /      \
       mul        mul
      /   \      /   \
    x0    x0   x1    x1
\end{BVerbatim}
\end{center}

This tree is called an \textbf{AST} (abstract syntax tree) --- a fancy name for
``the recipe, stored as data.'' Building it is \emph{all} that happens when you
write the formula.

\section{What the type checker did}

The \textbf{type checker} is a part of the Rust compiler. It runs \emph{before
your program ever executes}, when you build it. Its normal job is to catch
mismatches --- for instance, making sure you do not try to add a number to a
piece of text. Every value has a \textbf{type} (a label saying what kind of
thing it is: \code{Number}, \code{Text}, and so on), and the type checker makes
sure the labels are consistent.

Here is the trick. \textbf{A real number was hidden inside the type label.}

Normally a label just says ``this is a Number.'' Our labels say ``this is an
expression \emph{of degree 2}'' or ``\emph{of degree 0}.'' The \emph{degree} ---
the homogeneity weight, i.e.\ how the value scales when you scale the inputs ---
is baked into the type itself. And the labels carry arithmetic rules:

\begin{itemize}[leftmargin=1.4em]
  \item \code{var(0)} is labelled \textbf{degree 1}.
  \item \textbf{Multiplying} two expressions makes the result's degree the
  \textbf{sum} of the two: degree~1 times degree~1 gives degree~\textbf{2}.
  \item \textbf{Adding} requires both sides to have the \textbf{same} degree, and
  keeps it.
  \item \textbf{\code{sqrt}} \textbf{halves} the degree --- and is only allowed
  when the degree is even.
\end{itemize}

The striking part: because these rules live in the \emph{labels}, the type
checker is forced to \emph{do the degree arithmetic while it checks your code},
at compile time, before anything runs. Verifying that \code{var(0)*var(0)} is
well-typed literally computes $1+1=2$. Checking the \code{.sqrt()} checks that
$2$ is even and computes $2 \div 2 = 1$. The compiler has been turned into a
little calculator for degrees.

Two consequences, both at build time, not run time:

\begin{itemize}[leftmargin=1.4em]
  \item The formula above comes out labelled \textbf{degree 1}. Later, the
  evaluator reads that label to know how to rescale safely --- the degree was
  computed by the compiler, for free, with zero runtime cost.
  \item A formula whose degrees do not fit --- such as \code{x*x + 5}, which adds
  degree~2 to degree~0 --- violates the rule for \code{+} (both sides must
  match), and \textbf{the program refuses to compile.} Not a crash when it runs:
  it never builds, exactly as if you had tried to add a number to text.
\end{itemize}

\begin{keyidea}
A label like ``degree 2'' is normally static decoration. Here it is a
\emph{number the compiler adds, halves, and compares}. Most languages use types
only as coarse categories (``it's a Number''); this DSL uses the type as a
\emph{quantity the compiler does mathematics on}. The payoff: the compiler
proves a formula is scale-safe-able before it ever runs, and throws out the
formulas that are not --- with no runtime check and no way to forget it.
\end{keyidea}

\section{In what sense it is interpreted}

There are two ways to actually \emph{run} a language.

\begin{itemize}[leftmargin=1.4em]
  \item \textbf{Compiled.} Translate the whole recipe, once, into raw machine
  instructions the processor runs directly. Afterwards the recipe is gone; only
  the machine code remains. Fast.
  \item \textbf{Interpreted.} Keep the recipe around as data, and have a small
  program --- the \emph{interpreter} --- read it step by step, each time you want
  an answer, doing what each step says.
\end{itemize}

Our DSL is \textbf{interpreted}. Recall the tree of index cards. When you finally
ask for a number --- \code{norm.eval\_safe(\&[3.0, 4.0])} --- a function walks
that tree: \emph{``the top card says \code{sqrt}, so I first need the card below;
it says \code{add}, so I need its two children; this one says multiply
\code{x0} by \code{x0}, so look up $x_0 = 3$, multiply, get $9$}\,\ldots'' and so
on back up. That walking-the-tree-and-doing-what-it-says is interpretation. The
tree stays around; every evaluation re-walks it. The formula is \emph{never}
turned into dedicated machine code of its own.

\begin{keyidea}
One honest subtlety: the interpreter itself \emph{is} ordinary Rust, which Rust
compiles to machine code. So there is compiled code involved --- but it is the
code of the \emph{interpreter}, not of your \emph{formula}. Distinguish ``the
language the interpreter is written in'' (compiled Rust) from ``the little
language the interpreter runs'' (your formulas, interpreted). The
\code{compile\_safe} function, despite its name, does not compile anything: it
just prints the recipe as text.
\end{keyidea}

\section{The role of operator overloading}

The last piece: how is \code{var(0) * var(0)} even \emph{allowed} to mean ``build
a multiply card'' instead of ``multiply''?

\textbf{Operator overloading} means you get to decide what a symbol like \code{*}
or \code{+} does \emph{when its operands are your own type}. Normally \code{*}
means ``multiply two numbers,'' hardwired. Rust (like C++, Python, and others)
instead lets you say: ``when someone writes \code{a * b} and \code{a}, \code{b}
are my \code{Expr} type, here is what to do instead.'' What we told Rust to do
for \code{*} was, in plain words:

\begin{quote}
``Do not multiply. Build a new \emph{multiply} index card pointing at the two
operands. And set the result's degree label to the \emph{sum} of their degrees.''
\end{quote}

In code it is barely more than that:

\begin{center}
\begin{BVerbatim}
impl Mul for Expr {              // "here is what * means for my Expr type"
    fn mul(self, rhs) -> ... {
        Expr::new(Node::Mul(self, rhs))   // build a card; do NOT compute
    }
}
\end{BVerbatim}
\end{center}

Operator overloading is doing two jobs here:

\begin{enumerate}[leftmargin=1.6em]
  \item \textbf{Ergonomics.} It lets you write \code{x*x + y*y} with ordinary
  infix symbols instead of ugly nested calls like \code{add(mul(x,x),
  mul(y,y))}. The formula reads like mathematics. In most languages this nicer
  syntax is the \emph{only} reason overloading exists.
  \item \textbf{The hook.} Because we wrote the body of \code{*}, that body is
  exactly where both behaviours were inserted: ``build a card'' (so nothing
  computes yet) and ``add the degrees'' (the label arithmetic the type checker
  later enforces). Every overloaded operator is a little hook where the DSL's
  meaning is attached.
\end{enumerate}

\section{Where this pays off: the matrix operations}

The whole point of the machinery is that real linear-algebra routines should
read like their textbook definitions, with the safe rescaling generated rather
than typed by hand. There is one new idea needed to get there.

\subsection{One scaling knob becomes several}

The scalar DSL had a single degree because there was one way to scale the inputs:
multiply \emph{everything} by $2^{k}$. But an inner product
$\langle x, y\rangle = \sum_i x_i y_i$ has \emph{two independent} knobs --- you
can scale the $x$ vector and the $y$ vector separately --- and the result scales
by the product. So the label now carries a \emph{pair} of degrees,
$\langle D_x, D_y\rangle$. This is called \textbf{multigrading}.

\begin{itemize}[leftmargin=1.4em]
  \item \code{xvar()} (the $x$ vector) is labelled $\langle 1, 0\rangle$;
  \code{yvar()} is $\langle 0, 1\rangle$.
  \item Their elementwise product is $\langle 1, 1\rangle$ (the two labels add,
  componentwise).
  \item \code{sum(\ldots)} contracts a vector to a scalar and keeps the label.
  \item Adding requires both labels to match --- so \code{xvar() + yvar()}
  (that is $\langle 1,0\rangle + \langle 0,1\rangle$) is a \textbf{compile
  error}: you may not add an $x$-vector to a $y$-vector.
\end{itemize}

\subsection{Inner product: before and after}

Written by hand, the safe inner product carries the scaling openly --- find each
block's exponent, rescale every term, transport the sum back:

\begin{center}
\begin{BVerbatim}
let ex = block_exp(maxabs(x));
let ey = block_exp(maxabs(y));
let mut s = 0.0;
for i in 0..n { s += ldexp(x[i], -ex) * ldexp(y[i], -ey); }
ldexp(s, ex + ey)
\end{BVerbatim}
\end{center}

In the DSL the same operation is the naive formula. The label
$\langle 1,1\rangle$ is inferred by the compiler, and \code{eval\_safe} generates
the rescaling:

\begin{center}
\begin{BVerbatim}
let expr = sum(xvar() * yvar());   // bidegree <1,1>, inferred
expr.eval_safe(x, y)               // the block_exp/ldexp is generated
\end{BVerbatim}
\end{center}

\code{sum(x * y)} is simply the definition of an inner product.

\subsection{The other three}

\textbf{Matrix multiply} is ``each entry is the inner product of a row and a
column,'' delegating to the DSL inner product:

\begin{center}
\begin{BVerbatim}
for i in 0..rows {
    for j in 0..cols {
        c[i][j] = sum(xvar()*yvar()).eval_safe(row_i(A), col_j(B));
    }
}
\end{BVerbatim}
\end{center}

\textbf{Gram--Schmidt} and \textbf{eigenvalues (Jacobi)} are \emph{algorithms},
not single formulas --- they have loops and ``repeat until converged,'' which no
one expression can capture. They stay textbook by \emph{delegation}: every step
calls a safe primitive (\code{dot}, \code{nrm2}), so no scaling appears in them
either. Gram--Schmidt, in full:

\begin{center}
\begin{BVerbatim}
for j in 0..n {
    let mut w = v[j];
    for i in 0..j { let r = dot(q[i], w);  w -= r * q[i]; }  // remove projection
    let nrm = nrm2(w);
    if nrm > 0 { w /= nrm; }                                 // normalise
    q[j] = w;
}
\end{BVerbatim}
\end{center}

\begin{keyidea}
Count the scaling operations (\code{ldexp}) visible in the four routines:
\textbf{zero, zero, zero, and two.} Inner product and matrix multiply have none
--- the grading generated it. Gram--Schmidt has none --- it delegates to safe
primitives. Only Jacobi has two lines (rescale once at the top, transport the
eigenvalues once at the bottom), because an iterative algorithm cannot be a
single graded formula. Every routine now reads the way you would write it on a
whiteboard.
\end{keyidea}

There are, in fact, \emph{two} routes to ``looks naive but is safe,'' and the
matrix layer uses whichever fits. Inner product and matrix multiply are safe
by \textbf{grading}: you wrote the formula, a degree in the type generated the
rescale. Gram--Schmidt and Jacobi are safe by \textbf{composition}: you wrote
the algorithm, and it inherits safety from the primitives it calls --- which,
not coincidentally, is how LAPACK itself is organised.

\section{Putting the four together}

\begin{center}
\begin{BVerbatim}
You write:    (var(0)*var(0) + var(1)*var(1)).sqrt()
                        |
   operator overloading makes each  *  +  sqrt  build an index
   card instead of computing  -----------------------------+
                        |                                   |
                        v                                   v
   AT COMPILE TIME:                        AT RUN TIME:
   the type checker does the degree        eval_safe walks the
   arithmetic (1+1=2, 2/2=1), labels       card-tree and does the
   the result "degree 1", and REJECTS      actual numbers, reading
   any formula whose degrees don't fit     the "degree 1" label to
   (x*x + 5 won't even build)              rescale safely
\end{BVerbatim}
\end{center}

\begin{itemize}[leftmargin=1.4em]
  \item \textbf{Embedded DSL:} a little language made of ordinary Rust pieces
  that \emph{build a recipe} instead of computing.
  \item \textbf{Operator overloading:} the mechanism that lets \code{* + /} build
  recipe-cards (and carry the degree rules) instead of doing arithmetic.
  \item \textbf{Type checker:} a compile-time calculator that runs the degree
  arithmetic on the labels and throws out formulas that cannot be made safe ---
  all before the program runs.
  \item \textbf{Interpreted:} at run time, a function walks the recipe-tree step
  by step to get the number; your formula is never turned into its own machine
  code.
\end{itemize}

\end{document}
